% !TeX program = xelatex
% Сборка: два запуска xelatex TiAl_solvus_1000K.tex в одной папке.
% Все векторные рисунки и библиография включены непосредственно в исходник.
\documentclass[11pt,a4paper]{article}
\usepackage{fontspec}
\usepackage[russian]{babel}
\IfFontExistsTF{Cambria}{\setmainfont{Cambria}}{\setmainfont{DejaVu Serif}}
\IfFontExistsTF{Calibri}{\setsansfont{Calibri}}{\setsansfont{DejaVu Sans}}
\IfFontExistsTF{Consolas}{\setmonofont[Scale=MatchLowercase]{Consolas}}{\setmonofont{DejaVu Sans Mono}}
\usepackage[a4paper,left=23mm,right=21mm,top=21mm,bottom=20mm,headheight=15pt,headsep=7mm,footskip=10mm]{geometry}
\usepackage{amsmath,amssymb,booktabs,array,xcolor,titlesec,fancyhdr,caption,pgfplots,hyperref}
\pgfplotsset{compat=1.16}
\definecolor{ink}{HTML}{213542}
\definecolor{teal}{HTML}{006D70}
\definecolor{rust}{HTML}{A8492B}
\definecolor{muted}{HTML}{526673}
\definecolor{linegray}{HTML}{C5D0D3}
\hypersetup{unicode=true,colorlinks=true,linkcolor=teal,citecolor=teal,urlcolor=teal,
 pdftitle={Равновесные составы фаз α и α₂ в Ti–Al при 1000 К},pdfsubject={Кластерное разложение, свободная энергия и прямой расчёт на MTP}}
\urlstyle{same}
\captionsetup{font=small,labelfont={bf,color=teal},labelsep=period,singlelinecheck=false}
\titleformat{\section}{\sffamily\bfseries\Large\color{ink}}{\color{teal}\thesection}{.55em}{}
\titleformat{\subsection}{\sffamily\bfseries\large\color{ink}}{\thesubsection}{.55em}{}
\titlespacing*{\section}{0pt}{16pt plus 3pt}{6pt}
\titlespacing*{\subsection}{0pt}{10pt plus 2pt}{4pt}
\setlength{\parindent}{1.1em}
\setlength{\parskip}{2.5pt plus .7pt}
\setlength{\textfloatsep}{11pt plus 2pt minus 2pt}
\setlength{\intextsep}{10pt plus 2pt minus 2pt}
\setlength{\floatsep}{8pt plus 2pt}
\setlength{\emergencystretch}{2em}
\brokenpenalty=10000 \widowpenalty=8000 \clubpenalty=8000
\renewcommand{\arraystretch}{1.14}
\raggedbottom
\pagestyle{fancy}\fancyhf{}
\fancyhead[L]{\small\sffamily\color{muted}Ti--Al\enspace/\enspace Свободная энергия и фазовое равновесие}
\fancyhead[R]{\small\sffamily\color{muted}$1000$ К}
\fancyfoot[R]{\small\sffamily\color{muted}\thepage}
\renewcommand{\headrulewidth}{0.25pt}
\begin{document}
\thispagestyle{plain}
{\sffamily\small\color{teal}\bfseries НАУЧНЫЙ ОТЧЁТ\hfill 25 ноября 2026 г.\par}
\vspace{4mm}
{\sffamily\bfseries\fontsize{24}{28}\selectfont\color{ink}
Равновесные составы фаз \(\boldsymbol{\alpha}\) и \(\boldsymbol{\alpha_2}\)\\
в Ti--Al при 1000 К\par}
\vspace{2mm}
{\sffamily\large\color{muted}Прямая атомистическая оценка при нулевом давлении\\
и сравнение с кластерным разложением C46\par}
\vspace{4mm}\noindent{\color{teal}\rule{\linewidth}{1.1pt}}\par\vspace{1mm}

\noindent При свободном объёме и отношении осей \(c/a\) расчёт на потенциале
Qi и соавторов дал \textbf{14,98 ат.\% Al в \(\boldsymbol{\alpha}\)} и
\textbf{22,44 ат.\% Al в \(\boldsymbol{\alpha_2}\)}. Условные 95\%-е интервалы
составили \(14{,}92\)--\(15{,}06\) и \(22{,}33\)--\(22{,}55\) ат.\%.
Это равновесие двух кристаллических фаз в явно заданной области тепловых
смещений. Границы проверены на четырёх размерах, вплоть до 1400 атомов.
При изменении допустимых смещений и независимом повторении расчёта
измеренные составы менялись на величины порядка \(0{,}1\)--\(0{,}2\) процентного пункта.

\begin{center}
\begin{tikzpicture}
\begin{axis}[width=.91\linewidth,height=47mm,axis x line*=bottom,axis y line=none,
 xmin=14.0,xmax=23.15,ymin=.45,ymax=3.53,
 xtick={14,15,16,17,18,19,20,21,22,23},enlarge x limits=false,
 tick label style={font=\sffamily\footnotesize,text=muted},
 tick style={draw=muted},xlabel={Al, ат.\%},xlabel style={font=\sffamily\small},
 xmajorgrids,grid style={linegray!65},clip=false]
\addplot[teal,line width=1.5pt,mark=*,mark size=2.3pt] coordinates {(14.982,3) (22.445,3)};
\addplot[rust,line width=1.5pt,mark=square*,mark size=2.2pt] coordinates {(16.083,2) (22.159,2)};
\addplot[muted,line width=1.25pt,dashed,mark=triangle*,mark size=2.5pt] coordinates {(14.958,1) (21.958,1)};
\node[fill=white,text=teal,font=\sffamily\footnotesize,inner xsep=5pt] at (axis cs:18.60,3) {MTP, свободное \(c/a\)};
\node[fill=white,text=rust,font=\sffamily\footnotesize,inner xsep=5pt] at (axis cs:18.60,2) {MTP, \(c/a\) как у C46};
\node[fill=white,text=muted,font=\sffamily\footnotesize,inner xsep=5pt] at (axis cs:18.60,1) {C46, фиксированное \(c/a\)};
\node[above,text=ink,font=\small] at (axis cs:14.982,3.09) {\(\alpha\)};
\node[above,text=ink,font=\small] at (axis cs:22.445,3.09) {\(\alpha_2\)};
\end{axis}
\end{tikzpicture}
\captionof{figure}{Интервалы сосуществования при \(p=0\). Концы каждого
отрезка соответствуют равновесным составам двух фаз одной модели.
У прямого MTP \(N=1400\) и \(W<0{,}40\), у C46 \(N=5120\).}
\label{fig:tielines}
\end{center}

\noindent\textbf{Влияние формы ячейки.}
При одинаковом \(c/a\) прямой MTP даёт на \(1{,}12\) п.п. больше Al
в \(\alpha\) и на \(0{,}20\) п.п. больше в \(\alpha_2\), чем C46.
Освобождение отношения осей в атомистическом расчёте сдвигает границы
на \(-1{,}10\) и \(+0{,}29\) п.п. Таким образом, совпадение результатов
для \(\alpha\) при сравнении \(14{,}98\) с \(14{,}96\) возникает из-за
компенсации двух эффектов.

\noindent\textbf{Чувствительность к обучающей выборке.}
Независимая подгонка тех же 20 кластерных функций на структурах ATAT
даёт \(14{,}76\) и \(22{,}71\) ат.\% Al. При одинаковых механических
условиях обе подгонки предсказывают меньшую растворимость Al в \(\alpha\),
чем прямой MTP. Для \(\alpha_2\) прямое значение \(22{,}16\) лежит между
результатами двух подгонок. Здесь выбор обучающих структур меняет знак
отклонения, тогда как переход от среднего поля к Монте-Карло сильнее всего
влияет на границу \(\alpha\).

\noindent\textbf{Следующий расчёт.}
При тех же \(T,p,c/a\) целесообразно построить классический гармонический
вариант C46 и повторить определение обеих границ. Это даст явную оценку
влияния статистики ядер. Для измерения ошибки самого потенциала
нужна независимая разность свободных энергий MTP--DFT для обеих фаз.

\vfill
\noindent{\small\sffamily\color{muted}
Построение C46 — с.~\pageref{sec:ce-model}. Фазовые границы C46 — с.~\pageref{sec:ce-phases}.
Интегрирование в малых ячейках — с.~\pageref{sec:smallcells}.
Атомистический ансамбль — с.~\pageref{sec:ensemble}.
Составы и сходимость — с.~\pageref{sec:bulk-results}.
Дальнейшие расчёты — с.~\pageref{sec:outlook}.}
\clearpage
\section{Постановка задачи}
На богатой титаном стороне диаграммы Ti--Al фаза \(\alpha\) представляет
собой замещающий раствор Al в гексагональной плотноупакованной (ГПУ)
решётке. В фазе \(\alpha_2\) структуры D0$_{19}$ алюминий преимущественно
занимает одну из четырёх подрешёток. В идеальном \(\mathrm{Ti_3Al}\)
её заполнение полное, тогда как равновесная при конечной температуре
фаза допускает отклонение от стехиометрии.

Обозначим атомную долю алюминия через \(x=N_{\mathrm{Al}}/N\), а
свободную энергию Гиббса на атом фазы \(i\) через \(g_i(x,T,p)\).
Концы двухфазного интервала удовлетворяют условию общей касательной,
\begin{align}
  g'_\alpha(x_\alpha)&=g'_{\alpha_2}(x_{\alpha_2})=\Delta\mu, \\
  g_\alpha(x_\alpha)-x_\alpha\Delta\mu
    &=g_{\alpha_2}(x_{\alpha_2})-x_{\alpha_2}\Delta\mu ,
  \qquad \Delta\mu=\mu_{\rm Al}-\mu_{\rm Ti}.
\end{align}
Малый систематический сдвиг разности свободных энергий способен заметно
сместить касательную. Изменение свободной энергии одной декорации
нужно отличать от сдвига этой касательной после усреднения по всем
химическим заполнениям и атомным смещениям.

В проектной модели C46 энергии и колебательные вклады представлены
кластерным разложением. Исходные атомистические данные получены на
потенциале моментных тензоров (MTP) Qi и соавт.~\cite{qi2023}.
Воспроизведённая версия C46 использует авторские коэффициенты из комплекта
рукописи М.\,В.~Петрика, Е.\,О.~Хазиевой и Р.\,Е.~Рыльцева (версия v16).
Независимая подгонка на собственной выборке структур давала
при 1000 К \(14{,}764/22{,}707\) ат.\% Al. Сравнения ниже выполнены с
авторским набором \(14{,}958/21{,}958\) ат.\% Al и дополнены проверкой
чувствительности к выборке для подгонки.
В расчёте C46 химическое разупорядочение учитывалось методом
Монте-Карло, а колебания — в квантовом квазигармоническом приближении.
В настоящей работе использован тот же файл потенциала. Сначала фиксированные
декорации и полный перебор малых ячеек позволили проверить отдельные
вклады в свободную энергию. Затем для макроскопических составов мы
восстановили равновесные веса фаз из атомистических ансамблей без
кластерной аппроксимации энергии. С экспериментом прямые свободные
энергии Гиббса в этом расчёте не калибровались.

\paragraph{Три способа получить термодинамическую информацию.}
В решёточном MC интеграл \(\int\Delta\mu(x)\,dx\) восстанавливает
изотерму свободной энергии на атом, по которой проводится общая касательная.
В малых атомистических ячейках интеграл по параметру связи \(\lambda\)
измеряет работу перехода от гармонических сил к MTP при заданном
химическом заполнении. Для больших ячеек обе фазы отбирались в серии
перекрывающихся окон по параметру порядка \(Q\). Снятие смещающих полей
методом MBAR даёт относительные вероятности фаз, из которых следуют
равновесное поле и два состава. Различия этих процедур существенны при
переносе выводов от одной декорации к многоатомной системе.
